\subsection{Reference gaps and separately completed seeds}
For a reference policy $\pi$, define
\begin{equation}\label{eq:gaps}
\begin{gathered}
d_\pi(s,a)=q_\pi(s,\pi(s))-q_\pi(s,a),\qquad
\mathcal B_\pi=\{(s,a):d_\pi(s,a)>0\},\\
d_k(s,a)=q_k(s,\pi(s))-q_k(s,a).
\end{gathered}
\end{equation}
An action in $\mathcal B_\pi$ would worsen the reference action according to $q_\pi$. If these estimated gaps
are all positive, no greedy policy can select a bad action, so policy improvement applies. If
every $\mathcal B_\pi$ is empty, improvement applies between every pair of policies in both directions.
Their state values are equal; choosing a policy taking any prescribed action shows that every
action attains that same value. All policies are optimal with the same target, and Lemma~\ref{thm: bounded limit sets}
completes the theorem. Otherwise set
\begin{equation}\label{eq:constants}
d_* = \min_{\pi,\ (s,a)\in\mathcal B_\pi}d_\pi(s,a)>0,\qquad
\delta=d_*/4,\qquad h=g=\delta/2.
\end{equation}

Fix deterministic $K>B$ and restart $r$, and let $\zeta_r=\inf\{k\ge r:\|q_k\|_\infty>K\}$. For a
stopping time $t\ge r$, work on $\{t<\infty,t<\zeta_r\}$, fix $\pi=\pi_t$, and define
\begin{equation}\label{eq:stops}
\tau=\inf\{k>t:q_{\pi_k}\ne q_\pi\},\qquad \chi=\tau\wedge\zeta_r.
\end{equation}
These are stopping times. Return moments give finite second moments for every fixed-time
table by induction in \eqref{eq: update q with chi}. For $t\le k<\chi$, the action-value target is exactly $q_\pi$, even if the
selected policy label is different. Equal initial action probabilities give
\begin{align}
\mathbb E[d_{k+1}(s,a)-d_k(s,a)\mid\mathcal F'_k]
&=\alpha_k\sigma_k(s)(d_\pi(s,a)-d_k(s,a)),\label{eq:gapdrift}\\
\operatorname{Var}(d_{k+1}(s,a)-d_k(s,a)\mid\mathcal F'_k)
&\le V\alpha_k^2\sigma_k(s),\qquad V=2\{c_v+(B+K)^2\}.\label{eq:gapvariance}
\end{align}
Only the two coordinates forming the gap can change, their selection events are disjoint,
and each contributes at most $\alpha_k^2\sigma_k(s)\{c_v+(B+K)^2\}$ to the raw second moment. Since
$d_\pi(s,a)\ge4\delta$, the state-clock barrier applies up to and including $\chi$. All such moment estimates
are used before exit, not conditioned on future containment in the ball.

\paragraph{A favorable local update.}
Suppose the reference action is selected at $s$ and a current
bad gap is below $\delta$. Request one of the following action/return events:
\begin{enumerate}
\item If $q_\pi(s,\pi(s))-q_k(s,\pi(s))\ge d_*/2$, select $\pi(s)$ and require $G_{k+1}\ge q_\pi(s,\pi(s))-h$. The
reference coordinate increases by at least $3\delta\alpha_k/2$.
\item Otherwise select $a$ and require $G_{k+1}\le q_\pi(s,a)+h$. The identity
\[
q_k(s,a)-q_\pi(s,a)=q_k(s,\pi(s))-q_\pi(s,\pi(s))+d_\pi(s,a)-d_k(s,a)>\delta
\]
implies that the bad coordinate decreases by more than $g\alpha_k=\delta\alpha_k/2$.
\end{enumerate}
In either case no other gap against the reference action decreases, and the reference action
remains greedy. For centered $X$ with variance at most $c_v$, the one-sided Chebyshev inequality
gives
\[
\mathbb P(X\ge-h)\ge\frac{h^2}{c_v+h^2},\qquad
\mathbb P(X\le h)\ge\frac{h^2}{c_v+h^2}.
\]
For example, apply Markov's inequality to $(X+c_v/h)^2$ to bound $\mathbb P(X\ge h)$ and repeat for
$-X$; when $c_v=0$, $X=0$ almost surely. Conditional on selecting $s$, its initial action is
uniform, so every requested local update has conditional probability at least
\begin{equation}\label{eq:localprob}
p_0=\frac{h^2}{\max_s|\mathcal A(s)|\,(c_v+h^2)}>0.
\end{equation}
No probability of selecting the state enters this bound.

\begin{lemma}[Separately completed statewise seeds]\label{lem:seeds}
For full-inertia ties and any positive integer $N$, there is an adapted construction asking for at most $JN$
favorable local updates, starting at $t\ge k_\alpha$. Along successful requests, every unfinished state
retains its reference action, including at an endpoint reached by a favorable update. If state $s$
completes at time $T_s<\infty$, then
\begin{equation}\label{eq:seedmargin}
d_{T_s}(s,a)\ge\min\{\delta,gN\alpha_{T_s}/C_\alpha\}\qquad((s,a)\in\mathcal B_\pi).
\end{equation}
Each request, conditional on selecting its unfinished state, has probability at least $p_0$.
\end{lemma}
\begin{proof}
Order the bad actions at each state. Skip an action immediately if its gap is already
at least $\delta$; otherwise process it with at most $N$ favorable updates at that state, stopping
earlier if it reaches $\delta$. Process skips at $t$ and after every favorable update, before the next
sample. A state completes when its list is exhausted. Write $T_s$ for this time, using $T_s=t$ for
a state completed initially and $T_s=\infty$ if the construction stops before completion. These
are stopping times. Every selection of an unfinished state requests a local update; completed
states receive ordinary updates. Stop the construction at the first failed requirement or at $\chi$.

Under full inertia, raising the reference estimate or lowering a competitor retains the
reference action, since it remains greedy. At an unupdated state the table is unchanged, so
the action is also retained. Induction proves the assertion through every unfinished state's
wait.

The same retention conclusion also holds for fixed priorities: at time $t$ the reference action has highest priority among that state's
maximizers. Leaving the table unchanged, raising its estimate, or lowering a competitor
cannot create a new maximizer ahead of it in the priority order. Therefore that same action
continues to be selected at every unfinished state. Induct over all iterations, including those
updating other states and the final favorable update at completion. This proves the required
retention property directly; a fixed-priority rule is not being identified with full inertia on
arbitrary updates.

Every processed gap either reaches $\delta$ or receives $N$ increases of at least $g\alpha_k$. Later seed
updates at its state cannot lower it. Bounded future increases give $\alpha_k\ge\alpha_{T_s}/C_\alpha$
for every such update and proves \eqref{eq:seedmargin}. At most $J$ bad actions are processed, so there are
at most $JN$ requests, irrespective of waiting times. Their conditional probability bound is
\eqref{eq:localprob}.
\end{proof}

Completed states can be updated while other states wait. We therefore cannot condition
indiscriminately on all future successful behavior: the completed-state return laws must stay
unchanged. The following bounded change of measure conditions only the seed updates.

\begin{lemma}[Conditioning only the seed updates]\label{lem:measure}
There is an auxiliary law $\widehat{\mathbb P}\ll\mathbb P$
under which all active requirements of Lemma~\ref{lem:seeds} succeed. Conditional state selection is
unchanged. At completed states, and after the construction stops, the conditional action and
return law given the current history and selected state is unchanged. For every future event
$E$, on $\{t<\infty,t<\zeta_r\}$,
\begin{equation}\label{eq:transfer}
\mathbb P(E\mid\mathcal F'_t)\ge p_0^{JN}\widehat{\mathbb P}(E\mid\mathcal F'_t).
\end{equation}
\end{lemma}
\begin{proof}
Let $S_{k+1}$ be the state component of $I_{k+1}$. At an active unfinished state $s$, let
$p_k(s)$ be the conditional probability of its requirement given $\mathcal F'_k,S_{k+1}=s$. Then $p_k(s)\ge p_0$.
If the state has zero sampling probability, set $p_k(s)=p_0$; this value is never used on a
positive-probability draw.

Start $L_t=1$. At a requested update of $s$, multiply $L$ by the indicator of success divided
by $p_k(s)$; at other iterations multiply it by one. Every factor has conditional mean one given
$\mathcal F'_k,S_{k+1}$. Hence $(L_k)$ is a martingale. At most $JN$ factors differ from one, even though their
times need not be bounded, so $0\le L_k\le p_0^{-JN}$. Bounded convergence gives $\mathbb E[L_\infty\mid\mathcal F'_t]=1$.
Define $d\widehat{\mathbb P}=L_\infty\,d\mathbb P$. For random $t$, use likelihood one before $t$, on $\{t=\infty\}$, and off $\{t<\zeta_r\}$,
and combine the constructions on $\{t=j\}$ over integers $j$. The law agrees with $\mathbb P$ on $\mathcal F'_t$
where the construction is active.

A failure makes $L_\infty=0$, so active requirements succeed almost surely under the auxiliary
law. To check conditional laws, use $\mathbb E[L_\infty\mid\mathcal F'_{k+1}]=L_{k+1}$; future factors do not alter
the one-step likelihood. Its current factor has conditional mean one given the selected
state, preserving state selection. At a completed state the current factor is identically one,
so its conditional action and return law is unchanged as well. The statement holds on
auxiliary-almost-every history, where $L_k>0$. Absolute continuity preserves the almost-sure
algorithmic identities, including the full-inertia tie rule. Finally,
\[
\widehat{\mathbb P}(E\mid\mathcal F'_t)=\mathbb E[L_\infty\mathbf1_E\mid\mathcal F'_t]
\le p_0^{-JN}\mathbb P(E\mid\mathcal F'_t),
\]
which proves the bound. This auxiliary law is used only in the proof; the algorithm still
samples its original returns.
\end{proof}

